% 第6章 6.2：拉普拉斯的两条计算路线。单位权三点路径，f=(3,1,2)。
% (a) 逐点：(Lf)_2=(f2-f1)+(f2-f3)=(1-3)+(1-2)=-3；
% (b) 分步：B^T f=(2,-1)，W_E=I 不变，B(2,-1)=(2,-1-2,1)=(2,-3,1)。
\begin{tikzpicture}[x=1mm,y=1mm,>=Stealth,font=\figurefont\footnotesize,text=FigureInk,
  vertex/.style={circle,draw=FigureStroke,fill=FigureBlueFill,line width=.8pt,minimum size=6.5mm,inner sep=0pt},
  cur/.style={vertex,draw=FigureYellow,fill=FigureYellowFill},
  edge/.style={draw=FigureStroke,line width=.8pt},
  cell/.style={draw=FigureStroke,line width=.6pt,minimum width=7mm,minimum height=5mm,inner sep=0pt},
  vcell/.style={cell,fill=FigureBlueFill},
  ecell/.style={cell,fill=FigureYellowFill},
  flow/.style={->,draw=FigureStroke,line width=.75pt},
  note/.style={font=\FigureNoteFont\fontsize{7.5}{9}\selectfont,text=FigureMuted}]
% ---------- (a) ----------
\node[anchor=west] at (0,34) {(a) 逐点汇总};
\node[vertex] (a) at (4,18) {$1$};
\node[cur] (b) at (22,18) {$2$};
\node[vertex] (c) at (40,18) {$3$};
\draw[edge] (a)--(b)--(c);
\node[text=FigureBlue] at (4,25.5) {$3$};
\node[text=FigureBlue] at (22,25.5) {$1$};
\node[text=FigureBlue] at (40,25.5) {$2$};
\node[note] at (13,13) {$1-3$};
\node[note] at (31,13) {$1-2$};
\node at (22,5) {$(\bm L\bm f)_2=(-2)+(-1)=-3$};
% ---------- (b) ----------
\begin{scope}[xshift=52mm]
\node[anchor=west] at (0,34) {(b) 取边差—缩放—汇总};
\foreach \v [count=\i] in {3,1,2}{\node[vcell] (f\i) at (4,26-6*\i) {$\v$};}
\foreach \v [count=\i] in {2,-1}{\node[ecell] (d\i) at (21,23-6*\i) {$\v$};}
\foreach \v [count=\i] in {2,-1}{\node[ecell] (s\i) at (38,23-6*\i) {$\v$};}
\foreach \v [count=\i] in {2,-3,1}{\node[vcell] (y\i) at (55,26-6*\i) {$\v$};}
\draw[flow] (8.5,14)--(16.5,14) node[midway,above]{$\bm B^{\top}$};
\draw[flow] (25.5,14)--(33.5,14) node[midway,above]{$\bm W_E$};
\draw[flow] (42.5,14)--(50.5,14) node[midway,above]{$\bm B$};
\node[note] at (4,3) {$\bm f$};
\node[note] at (21,6) {边差};
\node[note] at (38,6) {乘边权};
\node[note] at (55,3) {$\bm L\bm f$};
\end{scope}
\end{tikzpicture}
